% ===========================================================================
% 1. Introduction
% ===========================================================================
\section{Introduction}
\label{sec:introduction}

The deployment of autonomous AI agents---systems capable of browsing the
web, executing code, sending messages, and making decisions on behalf of
users---has accelerated dramatically in recent years~\cite{yao2023react,
shinn2023reflexion, openai2024gpt4o}. Multi-agent orchestration
protocols such as Google's Agent-to-Agent (A2A)~\cite{google2025a2a}
further amplify both the capability and the attack surface of these
systems by enabling agents to delegate tasks, share context, and
coordinate actions across organizational boundaries.

Google DeepMind's seminal ``AI Agent Traps'' paper~\cite{yuntao2025agenttraps}
provides a rigorous taxonomy of six categories of environmental attacks
that exploit the operational context of AI agents rather than the models
themselves. The taxonomy spans \category{content injection} (hidden
instructions in rendered content), \category{semantic manipulation}
(social-engineering framing), \category{cognitive state corruption}
(RAG and knowledge-base poisoning), \category{behavioural control}
(deceptive UI elements), \category{systemic} attacks (multi-agent
cascade failures), and \category{human-in-the-loop} manipulation
(biased reporting to human reviewers). This contribution is
foundational---yet it remains \emph{purely theoretical}. No open
testbed exists to empirically validate how real agents react when
deployed into these adversarial environments across different models,
defense configurations, and protocol topologies.

\paragraph{The gap.}
Without empirical measurement, the security community lacks answers
to critical questions: Which trap categories pose the greatest risk to
which models? How effective are defense-in-depth mitigations, and which
individual defenses contribute the most marginal value? Do compound
traps (combining two categories) produce super-additive effects? Are
multi-agent cascade failures bounded or unbounded? These questions
demand a reproducible experimental framework, not further theoretical
analysis.

\paragraph{Our contribution.}
We present \textsc{agent-traps-lab}, an open-source testbed that
operationalizes all six DeepMind trap categories as 22 concrete
adversarial scenarios and evaluates them against GPT-4o-mini across
220 controlled experiments. Our contributions are:

\begin{enumerate}[nosep,leftmargin=*]
  \item \textbf{Testbed.} A modular, extensible framework implementing
    22 trap scenarios across all six categories, with a pluggable
    agent interface, mitigation registry, and automated experiment
    harness (\cref{sec:testbed}).

  \item \textbf{Experiment design.} A controlled comparison:
    22~scenarios $\times$ 2~conditions (baseline, hardened)
    $\times$ 5~repetitions = 220 experiment cells with
    Bonferroni-corrected non-parametric tests
    (\cref{sec:methodology}).

  \item \textbf{Empirical results.} Per-category trap success rates
    and detection rates with Cohen's $d$ effect sizes and Wilcoxon
    signed-rank tests
    (\cref{sec:results-ci}--\cref{sec:results-hitl}).

  \item \textbf{Defense evaluation.} A seven-module mitigation suite
    that reduces overall trap success from 45.5\% to 25.5\%, with
    category-specific analysis revealing both successes and regressions
    (\cref{sec:mitigations}).

  \item \textbf{Artifacts.} The full testbed, experiment data, and
    analysis scripts are publicly released for reproducibility
    (\cref{app:artifacts}).
\end{enumerate}

The remainder of this paper is organized as follows.
\Cref{sec:background} surveys the DeepMind taxonomy, the A2A protocol,
and multi-agent orchestration. \Cref{sec:testbed} describes our testbed
architecture. \Cref{sec:methodology} details the experiment design.
\Cref{sec:results-ci} through \cref{sec:results-hitl} present
per-category results. \Cref{sec:mitigations} reports overall mitigation
findings. \Cref{sec:discussion} discusses implications and limitations.
\Cref{sec:related} positions our work in the literature, and
\cref{sec:conclusion} concludes.
